\section{Method}\label{sec:method}

We present ColliForce-VLA, a framework that integrates differentiable workspace Signed Distance Fields (SDFs) into the forward pass of a pretrained Vision-Language-Action model to enable collision-aware action generation. The method requires no architectural modification to the VLA backbone; the default inference path uses only the base model with trained LoRA adapters. We describe the problem formulation (Section~\ref{sec:problem}), the differentiable SDF integration (Section~\ref{sec:sdf_integration}), the gradient-guided action refinement mechanism (Section~\ref{sec:gradient_refinement}), the training protocol (Section~\ref{sec:training}), and inference-time deployment (Section~\ref{sec:inference}).

\subsection{Problem Formulation}\label{sec:problem}

We consider a language-conditioned robotic manipulation setting where a Vision-Language-Action model $\pi_\theta$ generates an action trajectory $\{a_t\}_{t=1}^{H}$ with horizon $H = 50$, conditioned on visual observations $o$ and a language instruction $l$. Each action $a_t \in \mathbb{R}^7$ specifies the desired end-effector (EE) displacement in position ($\Delta x, \Delta y, \Delta z$), orientation ($\Delta q_1, \Delta q_2, \Delta q_3$), and gripper state.

We formalize collision avoidance through the workspace Signed Distance Field $\Phi: \mathbb{R}^3 \to \mathbb{R}$, where $\Phi(\mathbf{x}) > 0$ indicates free space, $\Phi(\mathbf{x}) = 0$ denotes the obstacle boundary, and $\Phi(\mathbf{x}) < 0$ indicates penetration. The objective is to learn a policy satisfying $\Phi(\mathrm{EE}(a_t)) > 0$ for all $t$ while preserving task success.

Our base model is $\pi_0$~\cite{black2024pi0}, a 3.5B-parameter flow-matching VLA comprising a PaliGemma vision-language encoder (Gemma 2B) and a Gemma 300M action expert. The action expert produces a sequence of action tokens through an output projection $W_{\mathrm{out}} \in \mathbb{R}^{1024 \times 32}$ that maps from the expert's hidden dimension to the action token space.

\subsection{Differentiable SDF Integration}\label{sec:sdf_integration}

We augment the VLA with two lightweight auxiliary modules that operate on the action expert's intermediate representations, together adding fewer than 660K parameters.

\paragraph{Feature extraction.} We capture the action expert's suffix output $\mathbf{z} \in \mathbb{R}^{B \times H \times 1024}$, the tensor immediately preceding $W_{\mathrm{out}}$, via a forward hook on \texttt{action\_out\_proj}. Mean-pooling along the action horizon dimension yields a compact representation:
\begin{equation}\label{eq:pooling}
\bar{\mathbf{z}} = \frac{1}{H} \sum_{t=1}^{H} \mathbf{z}_t \in \mathbb{R}^{1024}
\end{equation}
that summarizes the action expert's spatial understanding of the current scene.

\paragraph{SDF Decoder.} A 4-layer MLP $f_\mathrm{SDF}$ predicts the signed distance value at arbitrary 3D query points given the pooled features:
\begin{equation}\label{eq:sdf_decoder}
f_\mathrm{SDF}: (\bar{\mathbf{z}}, \mathbf{x}) \mapsto \hat{\Phi}(\mathbf{x}), \quad \bar{\mathbf{z}} \in \mathbb{R}^{1024},\; \mathbf{x} \in \mathbb{R}^3
\end{equation}
The architecture employs hidden dimension 256 with Softplus activation ($\beta = 5$) and a skip connection that re-injects the query coordinates at layer 2:
\begin{align}
\mathbf{h}_0 &= \sigma(\mathbf{W}_0 [\bar{\mathbf{z}}; \mathbf{x}] + \mathbf{b}_0), \quad \mathbf{W}_0 \in \mathbb{R}^{256 \times 1027} \label{eq:sdf_layer0} \\
\mathbf{h}_1 &= \sigma(\mathbf{W}_1 \mathbf{h}_0 + \mathbf{b}_1), \quad \mathbf{W}_1 \in \mathbb{R}^{256 \times 256} \\
\mathbf{h}_2 &= \sigma(\mathbf{W}_2 [\mathbf{h}_1; \mathbf{x}] + \mathbf{b}_2), \quad \mathbf{W}_2 \in \mathbb{R}^{256 \times 259} \\
\hat{\Phi}(\mathbf{x}) &= \mathbf{W}_3 \mathbf{h}_2 + \mathbf{b}_3, \quad \mathbf{W}_3 \in \mathbb{R}^{1 \times 256} \label{eq:sdf_output}
\end{align}
where $\sigma(\cdot) = \mathrm{Softplus}_{\beta=5}(\cdot)$. The total parameter count is ${\sim}$395K. Spatial gradients $\nabla_{\mathbf{x}} \hat{\Phi}(\mathbf{x})$ are computed via \texttt{torch.autograd.grad} with graph creation enabled during training, which is required for Eikonal regularization and the gradient-guided refinement in Section~\ref{sec:gradient_refinement}.

\paragraph{End-Effector (EE) Extractor.} A 2-layer MLP $g_\mathrm{EE}$ predicts the current EE position from the pooled features:
\begin{equation}\label{eq:ee_extractor}
g_\mathrm{EE}: \bar{\mathbf{z}} \mapsto \hat{\mathbf{p}}_\mathrm{EE} \in \mathbb{R}^3, \quad \mathrm{Linear}(1024, 256) \to \mathrm{ReLU} \to \mathrm{Linear}(256, 3)
\end{equation}
with ${\sim}$263K parameters. The EE prediction provides the query point for evaluating collision proximity, and its auxiliary loss encourages the action expert features to encode spatial position information.

\paragraph{Auxiliary losses.} The SDF decoder is trained with a clamped L1 loss and Eikonal regularization:
\begin{equation}\label{eq:loss_sdf}
\mathcal{L}_\mathrm{SDF} = \frac{1}{N} \sum_{i=1}^{N} \left| \mathrm{clamp}(\hat{\Phi}(\mathbf{x}_i), \pm\delta) - \mathrm{clamp}(\Phi^*(\mathbf{x}_i), \pm\delta) \right| + \lambda_\mathrm{eik} \frac{1}{N} \sum_{i=1}^{N} \left( \|\nabla_{\mathbf{x}} \hat{\Phi}(\mathbf{x}_i)\|_2 - 1 \right)^2
\end{equation}
where $\delta = 0.05$ is the clamping threshold (in meters), $\lambda_\mathrm{eik} = 0.1$, and $N = 64$ query points are sampled from the ground-truth SDF volume per step. The ground-truth SDF volume $\Phi^*(\mathbf{x})$ is precomputed from known obstacle meshes in the simulation workspace using mesh-to-SDF distance queries (via trimesh/pysdf) at a fixed spatial resolution \textbf{[TODO: confirm resolution from code]}. Clamping focuses learning on the near-surface region most relevant for collision detection; the Eikonal term enforces the unit-gradient property of valid distance fields. The EE extractor is trained with MSE loss:
\begin{equation}\label{eq:loss_ee}
\mathcal{L}_\mathrm{EE} = \|\hat{\mathbf{p}}_\mathrm{EE} - \mathbf{p}^*_\mathrm{EE}\|_2^2
\end{equation}
where $\mathbf{p}^*_\mathrm{EE} \in \mathbb{R}^3$ is the ground-truth end-effector position obtained from the robot's forward kinematics.

\subsection{SDF Gradient-Guided Action Refinement}\label{sec:gradient_refinement}

The core contribution of ColliForce-VLA is a gradient-based mechanism that steers the action expert's representations away from collision-prone regions during training. Unlike post-hoc action filtering or external safety layers, this approach modifies the feature space directly, enabling the model to internalize collision avoidance behavior into its LoRA parameters.

\paragraph{Refinement procedure.} Given the captured suffix output $\mathbf{z} \in \mathbb{R}^{B \times H \times 1024}$, the refinement proceeds in five steps:

\textbf{Step 1.} Predict EE position and evaluate the SDF:
\begin{equation}
\hat{\mathbf{p}}_\mathrm{EE} = g_\mathrm{EE}(\bar{\mathbf{z}}), \quad \hat{\Phi}(\hat{\mathbf{p}}_\mathrm{EE}) = f_\mathrm{SDF}(\bar{\mathbf{z}}, \hat{\mathbf{p}}_\mathrm{EE})
\end{equation}

\textbf{Step 2.} Identify near-collision samples within the batch:
\begin{equation}\label{eq:collision_set}
\mathcal{S} = \{i : \hat{\Phi}(\hat{\mathbf{p}}_{\mathrm{EE},i}) < m\}, \quad m = 0.01 \text{ (meters)}
\end{equation}

\textbf{Step 3.} For samples in $\mathcal{S}$, compute the feature-space SDF gradient:
\begin{equation}\label{eq:feature_grad}
\nabla_{\mathbf{z}} \hat{\Phi} = \frac{\partial \hat{\Phi}(\hat{\mathbf{p}}_\mathrm{EE})}{\partial \mathbf{z}}
\end{equation}
This gradient captures how changes in the action expert's internal representation affect the predicted collision distance. It is distinct from the spatial gradient $\nabla_{\mathbf{x}} \hat{\Phi}$ used in Eikonal regularization: $\nabla_{\mathbf{z}} \hat{\Phi}$ operates in the 1024-dimensional feature space and encodes directions that move the entire action representation toward safety.

\textbf{Step 4.} Perform gradient ascent on the SDF value (toward higher, safer distances):
\begin{equation}\label{eq:gradient_ascent}
\mathbf{z}' = \mathbf{z} + \alpha \cdot \mathrm{clip}(\nabla_{\mathbf{z}} \hat{\Phi}, \gamma) \odot \mathbf{m}_\mathcal{S}
\end{equation}
where $\alpha = 0.1$ is the step size, $\mathrm{clip}(\cdot, \gamma)$ denotes per-sample L2 norm clipping with threshold $\gamma = 1.0$ for gradient stability, and $\mathbf{m}_\mathcal{S}$ is a binary mask selecting only near-collision samples.

\textbf{Step 5.} Project both original and refined features through the action output projection:
\begin{equation}\label{eq:projection}
\mathbf{v}_t = W_\mathrm{out} \mathbf{z}_t, \quad \mathbf{v}'_t = W_\mathrm{out} \mathbf{z}'_t
\end{equation}

\paragraph{Refinement loss.} The model is trained to directly produce collision-corrected actions by minimizing the discrepancy between original and refined predictions on near-collision samples, summed over all $H$ action timesteps:
\begin{equation}\label{eq:loss_refine}
\mathcal{L}_\mathrm{refine} = \frac{1}{|\mathcal{S}|} \sum_{i \in \mathcal{S}} \sum_{t=1}^{H} \|\mathbf{v}_{t,i} - \mathrm{sg}(\mathbf{v}'_{t,i})\|_2^2
\end{equation}
where $\mathrm{sg}(\cdot)$ is the stop-gradient operator and the outer sum runs over batch samples in the collision set $\mathcal{S}$. The loss averages over samples but sums over timesteps, so the per-sample refinement signal scales with the trajectory length $H$; this ensures that the gradient magnitude is proportional to the cumulative deviation across the full action horizon rather than the per-step average, which we found empirically more effective for learning coherent trajectory-level corrections. The gradient flows only through $\mathbf{v}_t \to \mathbf{z} \to \theta_\mathrm{LoRA}$, training the LoRA adapters to anticipate and avoid collisions. The refined target $\mathbf{v}'_t$ is detached and serves as a teacher signal that encodes collision-aware behavior. Because the student (the LoRA-adapted VLA) learns to replicate the refined actions directly, the SDF gradient computation is not required at inference time; the collision-avoidance capability is absorbed into the model parameters during training.

\subsection{Training Protocol}\label{sec:training}

We adopt a two-phase training protocol that ensures SDF prediction accuracy before engaging gradient-guided refinement.

\paragraph{Phase 1: Auxiliary module warmup ($0 \to 5\mathrm{K}$ steps).} The entire backbone, including LoRA parameters, is frozen. Only the SDF decoder (${\sim}$395K params) and EE extractor (${\sim}$263K params) are trained on detached features ($\bar{\mathbf{z}}.\mathrm{detach}()$) to avoid destabilizing the pretrained backbone. This phase establishes an accurate SDF predictor anchored to the backbone's current feature distribution.

\paragraph{Phase 2: Joint training ($5\mathrm{K} \to 30\mathrm{K}$ steps).} LoRA parameters are unfrozen with a linear gradient warmup over 100 steps. We consider three training variants that differ only in their loss composition:

\textit{$\pi_0$-LoRA} (baseline): Standard LoRA fine-tuning with flow-matching action loss only:
\begin{equation}
\mathcal{L}_{\pi_0\text{-LoRA}} = \mathcal{L}_\mathrm{action}
\end{equation}

\textit{ColliForce-Aux}: Adds SDF and EE auxiliary losses without gradient-guided refinement:
\begin{equation}\label{eq:loss_aux}
\mathcal{L}_\mathrm{Aux} = \mathcal{L}_\mathrm{action} + \lambda_\mathrm{SDF} \cdot \mathcal{L}_\mathrm{SDF} + \lambda_\mathrm{EE} \cdot \mathcal{L}_\mathrm{EE}
\end{equation}
with $\lambda_\mathrm{SDF} = 0.1$ and $\lambda_\mathrm{EE} = 0.01$. Gradients flow from the auxiliary losses through the pooled features into the LoRA adapters, encouraging the action expert to develop features that encode workspace geometry.

\textit{ColliForce} (full method): Extends ColliForce-Aux with the gradient-guided refinement loss:
\begin{equation}\label{eq:loss_full}
\mathcal{L}_\mathrm{ColliForce} = \mathcal{L}_\mathrm{action} + \lambda_\mathrm{SDF} \cdot \mathcal{L}_\mathrm{SDF} + \lambda_\mathrm{EE} \cdot \mathcal{L}_\mathrm{EE} + \lambda_\mathrm{refine} \cdot \mathcal{L}_\mathrm{refine}
\end{equation}
with $\lambda_\mathrm{refine} = 0.1$. The refinement loss is computed only when near-collision samples exist in the batch ($|\mathcal{S}| > 0$), adding zero overhead for collision-free steps.

\paragraph{LoRA configuration.} Adapters are injected into both the PaliGemma encoder (rank 16, $\alpha = 16$) and the action expert (rank 32, $\alpha = 32$), targeting attention projections ($W_q, W_k, W_v, W_o$) and MLP layers ($W_\mathrm{gate}, W_\mathrm{up}, W_\mathrm{down}$). The total trainable LoRA parameter count is ${\sim}$33.5M. All variants use AdamW optimizer with cosine learning rate decay from $2.5 \times 10^{-5}$ to $2.5 \times 10^{-6}$, batch size 32, and train for 30K steps total.

\subsection{Inference}\label{sec:inference}

At inference, only the base model with trained LoRA adapters is active. The standard forward pass generates actions through iterative flow-matching denoising; neither the auxiliary modules nor gradient computation is required. The LoRA parameters have internalized collision avoidance behavior during training, producing trajectories that maintain clearance from obstacles without additional computation.

For deployment scenarios with out-of-distribution obstacle configurations, ColliForce-VLA supports an optional workspace-frame SDF gradient correction as a complementary safety net. If the predicted SDF at the next EE position falls below the safety margin ($\hat{\Phi}(\hat{\mathbf{p}}_\mathrm{next}) < m$), the position component of the action is corrected:
\begin{equation}\label{eq:inference_correction}
a'_t[:3] = a_t[:3] + \alpha \cdot \frac{\nabla_{\mathbf{x}} \hat{\Phi}(\hat{\mathbf{p}}_\mathrm{next})}{\|\nabla_{\mathbf{x}} \hat{\Phi}(\hat{\mathbf{p}}_\mathrm{next})\|_2}
\end{equation}
The L2-normalization bounds correction magnitude to $\alpha$ regardless of gradient scale. This mechanism introduces zero additional latency when $\hat{\Phi} \geq m$ and serves primarily as a safety net; the ColliForce-trained model has already learned to avoid collisions and rarely triggers inference-time correction.